\section{Motivation}
ALSAT-2B is an Algerian Earth-observation satellite providing multispectral imagery at 10~m spatial resolution. For many applications -- land-cover classification, building detection, change detection -- this resolution is insufficient to resolve fine-scale features. Super-resolution (SR) offers a post-processing path to enhance spatial detail without new hardware.

However, deep-learning SR methods are trained to minimize pixel-wise loss against high-resolution references. This objective can be satisfied by generating plausible-looking details that were not present in the original scene -- a phenomenon termed \textbf{hallucination}. Hallucinated details may improve PSNR while degrading the spectral/structural fidelity that downstream tasks depend on.

\section{Research Questions}
\begin{enumerate}
    \item[RQ1.] Do deep-learning SR methods improve fidelity metrics (PSNR, SSIM, SAM, ERGAS) over bicubic interpolation on ALSAT-2B imagery?
    \item[RQ2.] Do SR methods introduce hallucinations (spectral, structural, textural, semantic), and at what rate?
    \item[RQ3.] Does SR reduce or increase hallucination compared to the non-learned bicubic baseline?
\end{enumerate}

\section{Contributions}
\begin{enumerate}
    \item A hallucination taxonomy and detection framework with four categories (spectral, structural, textural, semantic) and configurable thresholds.
    \item A simulated ALSAT-2B $\times$4 dataset derived from Sentinel-2 (8 AOIs, 4 bands, 10{,}697 train+val tiles).
    \item Empirical evidence that EDSR-small improves fidelity (+3.03~dB PSNR) AND reduces hallucination rate (6.89\% vs 13.03\%) compared to bicubic.
    \item An open, reproducible pipeline (code, configs, scripts) for SR evaluation on satellite imagery, including a Python GUI application.
\end{enumerate}

\section{Thesis Structure}
Chapter~2 reviews related work. Chapter~3 describes the methodology. Chapter~4 presents baseline implementation. Chapter~5 reports SR experiments. Chapter~6 presents hallucination-aware evaluation. Chapter~7 discusses results. Chapter~8 concludes.
